\chapter{Measurement Validity: Validating the Ruler}
\label{ch:measurement-validity}

\section{Why This Chapter Exists}
\label{sec:validity-why}

A measurement claim rests on two things: the system measured, and the instrument that measured
it. The document-understanding literature spends most of its effort on the first. This chapter
is about the second, and it exists because on this task the instrument turned out to be the
larger source of error.

The claim is specific and it is not a generality about careful work. Over the course of this
project, a sequence of defects was found in the evaluation apparatus --- in the reference data,
in the comparison rules, in the diagnostic aggregation, and in the field specification given to
the model. Each was found by opening individual documents, not by looking at aggregates.
Each had been invisible in every number reported up to the point it was found. Several of them
moved reported scores by more than any modelling change made in this project. One of them, had
it gone unnoticed, would have caused this thesis to report the opposite of its actual
conclusion.

That is a result about method, and it is reported as one. It also carries a warning that
generalises past this system: \textbf{a document-extraction score computed by naive comparison
against canonical reference data is systematically pessimistic, and the size of the bias is not
small}. A reader who takes one thing from this thesis should take that.

The chapter proceeds by stating the discipline that separates repairing an instrument from
inflating a score (\S\ref{sec:validity-discipline}), then the four classes of defect found
(\S\ref{sec:validity-reference}--\S\ref{sec:validity-specification}), then what correcting them
was worth (\S\ref{sec:validity-effect}), then the near-miss
(\S\ref{sec:validity-confound}), and finally what remains broken and is recorded instead of
repaired (\S\ref{sec:validity-open}). Figure~\ref{fig:defect-chronology} gives the whole
sequence in one view, with each defect class pointing at the evidence that carries its
effect.

\begin{figure}[htbp]
  \centering
  %% The four defect classes of Chapter 6 (plus the self-inflicted fix and the
%% near-miss), arranged as the audit sequence they were found in, each pointing at
%% the evidence that carries its effect. Deliberately qualitative: every magnitude
%% lives in a generated table, and this figure only routes the reader to them.
%%
%% Hand-authored TikZ (not generated): it depicts the structure of the chapter's
%% argument -- found by opening documents, fixed under the four-rule discipline,
%% measured on frozen generations -- not a measurement itself.
%%
%% Input from a chapter with:  %% The four defect classes of Chapter 6 (plus the self-inflicted fix and the
%% near-miss), arranged as the audit sequence they were found in, each pointing at
%% the evidence that carries its effect. Deliberately qualitative: every magnitude
%% lives in a generated table, and this figure only routes the reader to them.
%%
%% Hand-authored TikZ (not generated): it depicts the structure of the chapter's
%% argument -- found by opening documents, fixed under the four-rule discipline,
%% measured on frozen generations -- not a measurement itself.
%%
%% Input from a chapter with:  %% The four defect classes of Chapter 6 (plus the self-inflicted fix and the
%% near-miss), arranged as the audit sequence they were found in, each pointing at
%% the evidence that carries its effect. Deliberately qualitative: every magnitude
%% lives in a generated table, and this figure only routes the reader to them.
%%
%% Hand-authored TikZ (not generated): it depicts the structure of the chapter's
%% argument -- found by opening documents, fixed under the four-rule discipline,
%% measured on frozen generations -- not a measurement itself.
%%
%% Input from a chapter with:  \input{figures/defect-chronology.tex}

\begin{tikzpicture}[
    font=\small,
    class/.style={
      draw=horusaccent, line width=0.8pt, rounded corners=2pt,
      fill=horuswash, text width=5.15cm, align=left, inner sep=5pt,
    },
    selfclass/.style={
      draw=horuswarn, line width=0.8pt, rounded corners=2pt,
      fill=white, text width=5.15cm, align=left, inner sep=5pt,
    },
    spine/.style={line width=1.1pt, draw=horusmuted},
    tick/.style={fill=horusaccent, circle, inner sep=1.6pt},
    lbl/.style={font=\scriptsize, text=horusink},
  ]

  %% Vertical spine: the audit sequence. Two columns of class boxes alternate sides
  %% so six boxes fit one column width.
  \draw[spine] (0,0.6) -- (0,-10.9);
  \node[lbl, anchor=south] at (0,0.7) {\textbf{found by opening individual documents
    --- never visible in an aggregate}};

  \node[tick] at (0,-0.6) {};
  \node[class, anchor=east] (c1) at (-0.45,-0.9) {%
    \textbf{Class 1 --- reference data wrong or ill-posed}\\[1pt]
    \scriptsize omitted optional totals, undefined document-level scalars, incomplete
    code maps, unlocatable answer-key values\\[2pt]
    \color{horuswarn}\itshape effect: Table~\ref{tab:heldout-headline} protocol note%
  };

  \node[tick] at (0,-2.4) {};
  \node[class, anchor=west] (c2) at (0.45,-2.7) {%
    \textbf{Class 2 --- comparison rule punished correct answers}\\[1pt]
    \scriptsize labels on values, controlled vocabularies the page never prints,
    structural prefixes, sign conventions\\[2pt]
    \color{horuswarn}\itshape effect: Table~\ref{tab:reader-lineage}, rows 2 vs.\ 4%
  };

  \node[tick] at (0,-4.2) {};
  \node[class, anchor=east] (c3) at (-0.45,-4.5) {%
    \textbf{Class 3 --- the apparatus itself was broken}\\[1pt]
    \scriptsize contaminated per-field diagnostic, unfaithful perfect-text renderer,
    blank-page-blind reading proxy\\[2pt]
    \color{horuswarn}\itshape effect: Table~\ref{tab:oracle-renderer-correction},
    Figure~\ref{fig:ruler-correction}%
  };

  \node[tick] at (0,-6.0) {};
  \node[class, anchor=west] (c4) at (0.45,-6.3) {%
    \textbf{Class 4 --- specification not grounded in the data}\\[1pt]
    \scriptsize 34 of 63 label variants never occur in the corpus; four reference
    values leaked as ``examples''\\[2pt]
    \color{horuswarn}\itshape now gated by a test and a corpus-backed audit%
  };

  \node[tick, fill=horuswarn] at (0,-7.8) {};
  \node[selfclass, anchor=east] (c5) at (-0.45,-8.1) {%
    \textbf{A correction that was itself a defect}\\[1pt]
    \scriptsize the sign-convention fix was applied one-sidedly and scored a correct
    negative answer as a miss --- the symmetry rule exists because this happened\\[2pt]
    \color{horuswarn}\itshape \S\ref{sec:validity-selfinflicted}%
  };

  \node[tick, fill=horuswarn] at (0,-9.6) {};
  \node[selfclass, anchor=west] (c6) at (0.45,-9.9) {%
    \textbf{The near-miss --- precision confound}\\[1pt]
    \scriptsize full-precision adapter against a 4-bit baseline: two correct numbers,
    one invalid comparison, caught only by a pre-registered rule\\[2pt]
    \color{horuswarn}\itshape effect: Table~\ref{tab:precision-confound}%
  };

  \node[
    draw=horusmuted, line width=0.5pt, rounded corners=2pt, dotted,
    fill=white, text=horusink, font=\scriptsize, align=center,
    text width=9.6cm, inner sep=4pt, anchor=north,
  ] at (0,-11.3) {Every correction was measured on \textbf{frozen generations} ---
    no model ran, so each movement is attributable to the instrument alone --- under
    the four-rule discipline of \S\ref{sec:validity-discipline}.};

\end{tikzpicture}


\begin{tikzpicture}[
    font=\small,
    class/.style={
      draw=horusaccent, line width=0.8pt, rounded corners=2pt,
      fill=horuswash, text width=5.15cm, align=left, inner sep=5pt,
    },
    selfclass/.style={
      draw=horuswarn, line width=0.8pt, rounded corners=2pt,
      fill=white, text width=5.15cm, align=left, inner sep=5pt,
    },
    spine/.style={line width=1.1pt, draw=horusmuted},
    tick/.style={fill=horusaccent, circle, inner sep=1.6pt},
    lbl/.style={font=\scriptsize, text=horusink},
  ]

  %% Vertical spine: the audit sequence. Two columns of class boxes alternate sides
  %% so six boxes fit one column width.
  \draw[spine] (0,0.6) -- (0,-10.9);
  \node[lbl, anchor=south] at (0,0.7) {\textbf{found by opening individual documents
    --- never visible in an aggregate}};

  \node[tick] at (0,-0.6) {};
  \node[class, anchor=east] (c1) at (-0.45,-0.9) {%
    \textbf{Class 1 --- reference data wrong or ill-posed}\\[1pt]
    \scriptsize omitted optional totals, undefined document-level scalars, incomplete
    code maps, unlocatable answer-key values\\[2pt]
    \color{horuswarn}\itshape effect: Table~\ref{tab:heldout-headline} protocol note%
  };

  \node[tick] at (0,-2.4) {};
  \node[class, anchor=west] (c2) at (0.45,-2.7) {%
    \textbf{Class 2 --- comparison rule punished correct answers}\\[1pt]
    \scriptsize labels on values, controlled vocabularies the page never prints,
    structural prefixes, sign conventions\\[2pt]
    \color{horuswarn}\itshape effect: Table~\ref{tab:reader-lineage}, rows 2 vs.\ 4%
  };

  \node[tick] at (0,-4.2) {};
  \node[class, anchor=east] (c3) at (-0.45,-4.5) {%
    \textbf{Class 3 --- the apparatus itself was broken}\\[1pt]
    \scriptsize contaminated per-field diagnostic, unfaithful perfect-text renderer,
    blank-page-blind reading proxy\\[2pt]
    \color{horuswarn}\itshape effect: Table~\ref{tab:oracle-renderer-correction},
    Figure~\ref{fig:ruler-correction}%
  };

  \node[tick] at (0,-6.0) {};
  \node[class, anchor=west] (c4) at (0.45,-6.3) {%
    \textbf{Class 4 --- specification not grounded in the data}\\[1pt]
    \scriptsize 34 of 63 label variants never occur in the corpus; four reference
    values leaked as ``examples''\\[2pt]
    \color{horuswarn}\itshape now gated by a test and a corpus-backed audit%
  };

  \node[tick, fill=horuswarn] at (0,-7.8) {};
  \node[selfclass, anchor=east] (c5) at (-0.45,-8.1) {%
    \textbf{A correction that was itself a defect}\\[1pt]
    \scriptsize the sign-convention fix was applied one-sidedly and scored a correct
    negative answer as a miss --- the symmetry rule exists because this happened\\[2pt]
    \color{horuswarn}\itshape \S\ref{sec:validity-selfinflicted}%
  };

  \node[tick, fill=horuswarn] at (0,-9.6) {};
  \node[selfclass, anchor=west] (c6) at (0.45,-9.9) {%
    \textbf{The near-miss --- precision confound}\\[1pt]
    \scriptsize full-precision adapter against a 4-bit baseline: two correct numbers,
    one invalid comparison, caught only by a pre-registered rule\\[2pt]
    \color{horuswarn}\itshape effect: Table~\ref{tab:precision-confound}%
  };

  \node[
    draw=horusmuted, line width=0.5pt, rounded corners=2pt, dotted,
    fill=white, text=horusink, font=\scriptsize, align=center,
    text width=9.6cm, inner sep=4pt, anchor=north,
  ] at (0,-11.3) {Every correction was measured on \textbf{frozen generations} ---
    no model ran, so each movement is attributable to the instrument alone --- under
    the four-rule discipline of \S\ref{sec:validity-discipline}.};

\end{tikzpicture}


\begin{tikzpicture}[
    font=\small,
    class/.style={
      draw=horusaccent, line width=0.8pt, rounded corners=2pt,
      fill=horuswash, text width=5.15cm, align=left, inner sep=5pt,
    },
    selfclass/.style={
      draw=horuswarn, line width=0.8pt, rounded corners=2pt,
      fill=white, text width=5.15cm, align=left, inner sep=5pt,
    },
    spine/.style={line width=1.1pt, draw=horusmuted},
    tick/.style={fill=horusaccent, circle, inner sep=1.6pt},
    lbl/.style={font=\scriptsize, text=horusink},
  ]

  %% Vertical spine: the audit sequence. Two columns of class boxes alternate sides
  %% so six boxes fit one column width.
  \draw[spine] (0,0.6) -- (0,-10.9);
  \node[lbl, anchor=south] at (0,0.7) {\textbf{found by opening individual documents
    --- never visible in an aggregate}};

  \node[tick] at (0,-0.6) {};
  \node[class, anchor=east] (c1) at (-0.45,-0.9) {%
    \textbf{Class 1 --- reference data wrong or ill-posed}\\[1pt]
    \scriptsize omitted optional totals, undefined document-level scalars, incomplete
    code maps, unlocatable answer-key values\\[2pt]
    \color{horuswarn}\itshape effect: Table~\ref{tab:heldout-headline} protocol note%
  };

  \node[tick] at (0,-2.4) {};
  \node[class, anchor=west] (c2) at (0.45,-2.7) {%
    \textbf{Class 2 --- comparison rule punished correct answers}\\[1pt]
    \scriptsize labels on values, controlled vocabularies the page never prints,
    structural prefixes, sign conventions\\[2pt]
    \color{horuswarn}\itshape effect: Table~\ref{tab:reader-lineage}, rows 2 vs.\ 4%
  };

  \node[tick] at (0,-4.2) {};
  \node[class, anchor=east] (c3) at (-0.45,-4.5) {%
    \textbf{Class 3 --- the apparatus itself was broken}\\[1pt]
    \scriptsize contaminated per-field diagnostic, unfaithful perfect-text renderer,
    blank-page-blind reading proxy\\[2pt]
    \color{horuswarn}\itshape effect: Table~\ref{tab:oracle-renderer-correction},
    Figure~\ref{fig:ruler-correction}%
  };

  \node[tick] at (0,-6.0) {};
  \node[class, anchor=west] (c4) at (0.45,-6.3) {%
    \textbf{Class 4 --- specification not grounded in the data}\\[1pt]
    \scriptsize 34 of 63 label variants never occur in the corpus; four reference
    values leaked as ``examples''\\[2pt]
    \color{horuswarn}\itshape now gated by a test and a corpus-backed audit%
  };

  \node[tick, fill=horuswarn] at (0,-7.8) {};
  \node[selfclass, anchor=east] (c5) at (-0.45,-8.1) {%
    \textbf{A correction that was itself a defect}\\[1pt]
    \scriptsize the sign-convention fix was applied one-sidedly and scored a correct
    negative answer as a miss --- the symmetry rule exists because this happened\\[2pt]
    \color{horuswarn}\itshape \S\ref{sec:validity-selfinflicted}%
  };

  \node[tick, fill=horuswarn] at (0,-9.6) {};
  \node[selfclass, anchor=west] (c6) at (0.45,-9.9) {%
    \textbf{The near-miss --- precision confound}\\[1pt]
    \scriptsize full-precision adapter against a 4-bit baseline: two correct numbers,
    one invalid comparison, caught only by a pre-registered rule\\[2pt]
    \color{horuswarn}\itshape effect: Table~\ref{tab:precision-confound}%
  };

  \node[
    draw=horusmuted, line width=0.5pt, rounded corners=2pt, dotted,
    fill=white, text=horusink, font=\scriptsize, align=center,
    text width=9.6cm, inner sep=4pt, anchor=north,
  ] at (0,-11.3) {Every correction was measured on \textbf{frozen generations} ---
    no model ran, so each movement is attributable to the instrument alone --- under
    the four-rule discipline of \S\ref{sec:validity-discipline}.};

\end{tikzpicture}

  \caption[The defect classes and their evidence]{The defect classes of this chapter, in
  the order the audit surfaced them. Every one was found by opening individual documents,
  repaired under the four-rule discipline of \S\ref{sec:validity-discipline}, and measured
  on frozen generations; the two red entries are the project's own errors, reported
  because they are the strongest argument for the rules.}
  \label{fig:defect-chronology}
\end{figure}

\section{The Discipline: Repair Versus Inflation}
\label{sec:validity-discipline}

Every correction described in this chapter raises a score. That is exactly the property that
makes the activity dangerous, because a sufficiently motivated experimenter can improve any
result by adjusting the scoring rules until it looks better. Four rules separate the two
activities, and they were fixed before the corrections began.

\textbf{A correction may only change how a representation is compared, never whether a wrong
answer counts as right.} Stripping a label from a predicted invoice number so that
\emph{Invoice no.\ 471102} matches the reference \emph{471102} is a representation change: the
model read the number correctly. Accepting \emph{471103} would not be. Every correction in this
chapter is of the first kind, and each one carries a regression test asserting that a
genuinely wrong value still fails.

\textbf{Any rule applied to a predicted value is applied identically to the reference value.}
Asymmetric normalisation is the specific mechanism by which a fairness fix becomes a scoring
error, and \S\ref{sec:validity-selfinflicted} reports the occasion on which this project
violated its own rule and what it cost.

\textbf{The effect of a correction is measured on frozen model outputs.} Generations are
archived, so a scoring change can be applied to identical stored predictions and the entire
difference attributed to the instrument. Had the outputs been regenerated, a representation fix
and a model change would have been inseparable. This is the single most useful practice
described in this chapter, and it is cheap: it requires only that generations be written to
disk.

\textbf{An exclusion requires a criterion declared in advance and applied by a rule.} A cell
may be removed from scoring only on a stated, general warrant --- never case by case after
observing whether the model happened to get it right. Where a criterion could not be
generalised safely, it was deliberately confined to the one field the evidence supported, and
\S\ref{sec:validity-reference} says why.

\section{Class One: The Reference Data Was Wrong or Ill-Posed}
\label{sec:validity-reference}

The reference for the synthetic corpus is extracted from each document's own embedded data, so
it cannot disagree with itself. It can still be an unfair standard, in three distinct ways.

\textbf{Optional values the reference omits and the page implies.} The invoice standard makes
several totals optional. A document with no discounts may omit the discount total entirely
rather than record it as zero. A model reading the page may reasonably emit zero. Scoring that
as a fabrication punishes a correct reading of an absent quantity. The correction treats an
omitted optional total and an explicit zero as the same claim --- in both directions, so a
model that emits a non-zero value where the reference omits one still fails.

\textbf{Document-level scalars that are undefined for some documents.} The schema carries a
single tax rate, inherited from the era when an invoice had one. A document carrying two
different rates has no single rate, so no answer is correct and the field is ill-posed for that
document. It is excluded, not scored --- and separately, an invoice whose only rate is
zero per cent presents the same problem in a different guise, because a zero rate is
indistinguishable from an absent one by value alone. Both exclusions are rule-driven on a
declared criterion.

\textbf{Reference data that is simply incomplete.} The document-type code map initially covered
the common invoice codes and not the full family the standard defines, so documents typed with
a less common code were scored against nothing. A construction invoice and an insurer's invoice
carry their own codes; the page prints the German phrase, the reference stores the numeric code,
and the map between them has to be complete or the field fails for reasons that have nothing to
do with reading. Similarly, payment terms are sometimes carried only as free text, so a
reference built solely from structured fields records an absence the page contradicts.

\textbf{The held-out corpus needed a fourth kind of correction}, and this one deserves care
because it is the most easily abused. Eight cells in the signed-off answer key carry values
that \emph{no} reading channel could locate anywhere in the page text --- not the vision model,
not the commercial service, not the text-layer parser, not the author reading the image. A
value that is not on the page cannot be extracted from the page by any vision system, so
scoring it as a miss measures the answer key, not the system.

The temptation is to generalise that criterion. It was deliberately not generalised, for a
reason worth stating: the criterion is confounded across fields. One amount is unlocatable on a
dozen documents and yet correctly recovered on eleven of them, because it is fixed
arithmetically by other fields without being read. Keying the exclusion on provenance class alone
would therefore have discarded correct answers. The exclusion was applied to exactly one field,
on evidence specific to it: every present cell of that field rests on an author decision with
no printed proof and no cross-channel agreement, the standard derives it from no other field,
and the model recovers none of them while scoring perfectly on the same field given perfect
text. The effect is reported in Table~\ref{tab:heldout-headline}: recall improves, the false-
negative count falls, and \textbf{the true-positive and false-positive counts do not move at
all}. Not one correct answer was discarded to obtain it.

\section{Class Two: The Comparison Rule Punished Correct Answers}
\label{sec:validity-comparison}

The second class is the gap between an exact reference and a fair comparison. The reference
stores canonical machine values; the page prints what a human reads. Four examples stand for
the class.

\textbf{Values printed with their label.} An invoice prints \emph{Invoice no.\ 471102}. A model
transcribing and extracting faithfully may return the whole phrase. The number is right; the
string differs. The correction strips a recognised leading label from the predicted side before
comparison.

\textbf{Controlled vocabularies the page never prints.} The standard encodes tax categories as
letters and payment methods as numeric codes. No invoice prints \emph{category S}; it prints a
German phrase describing the tax treatment. A model that reads the phrase and reports it has
read correctly, and a comparison against the code letter scores it wrong. The correction maps
the printed forms onto the coded ones.

\textbf{Identifiers carrying a structural prefix.} A supplier's article number may be printed
with a global trade item number prefix that the reference stores separately. Comparing the two
as strings fails on a correct reading.

\textbf{Sign conventions.} Discount and prepayment totals are stored unsigned in most
documents and printed with a minus sign. A model that emits the printed form is right and
scores as a miss.

None of these is exotic. All four are the ordinary consequence of measuring a
page-reading system against machine-readable reference data, which is why the class matters
beyond this project: any evaluation constructed the same way inherits the same bias, and
inherits it silently.

\section{Class Three: The Apparatus Itself Was Broken}
\label{sec:validity-apparatus}

The third class is the most uncomfortable, because these are defects in the measuring code
rather than in the difficulty of the task.

\textbf{The per-field diagnostic was contaminated by non-signal outcomes.} A correct absence
and an excluded cell were both being accumulated into each field's reported mean --- the first
contributing a perfect score, the second a zero. Neither carries information about how well a
field was read. The consequence is that a field's reported number became partly a function of
\emph{how often that field happens to be absent}: a field that is usually absent and correctly
predicted absent drifted upward toward a perfect score, while a field with many legitimate
exclusions drifted toward zero.

The practical damage is worth stating precisely, because it is the clearest example in this
project of a diagnostic defect with real consequences. The plan at the time was to aim the
adaptation experiment at the weakest fields. The contaminated ranking put the document-level tax
rate near the bottom and the rounding amount at a perfect score. \textbf{Both were artefacts,
and the ranking was wrong in both directions.} An adaptation study aimed using that ranking
would have targeted a field that was fine and ignored one that was not.

The headline metrics were never affected --- they derive from true-positive, false-positive and
false-negative counts, which were always computed correctly. This was a diagnostic defect. But
the diagnostic is what aims the work, so a defect there propagates into every decision about
what to fix next. The repair introduced a single shared predicate for signal-bearing outcomes,
used at all three aggregation sites, and a field with no signal-bearing outcomes is now omitted
instead of reported as perfect. Reporting a perfect score there says \emph{always right} when
it means \emph{never asked}.

\textbf{The perfect-text instrument was unfaithful.} Instrument three of
\S\ref{sec:method-attribution} renders text from the reference data and feeds it to the
structuring stage. If that renderer omits or misstates something the reference contains, the
resulting ceiling is wrong --- and wrong in an unpredictable direction, not merely
optimistic, because the structurer may be penalised for a value the renderer never showed it.
Five defects were found in it. Two fields that had been scoring zero on perfect text turned out
to be scoring zero because the renderer never rendered them; after repair they score near the
top of the range. Table~\ref{tab:oracle-renderer-correction} reports the whole effect on
identical archived generations.

% GENERATED FILE -- DO NOT EDIT BY HAND.
% Regenerate with: make thesis-assets  (scripts/thesis_assets.py)
% Editing this file by hand breaks the guarantee that every number in the
% manuscript is derived from committed evidence (ADR-055).

% Source(s): data/finetune/eval-oracle-tier1-val.json; data/finetune/eval-oracle-adr059-val.json

\begin{table}[htbp]
  \centering
  \caption{Repairing the perfect-text renderer. Measured on identical archived generations, so the whole difference is attributable to the measuring instrument.}
  \label{tab:oracle-renderer-correction}
  \begin{tabular}{lrrr}
    \toprule
    Metric & Before & After & $\Delta$ \\
    \midrule
    Overall F$_1$ & 0.9676 & 0.9719 & +0.0043 \\
    Flat F$_1$ & 0.9641 & 0.9743 & +0.0102 \\
    Presence-conditional F$_1$ & 0.9717 & 0.9812 & +0.0095 \\
    Spurious rate & 0.0176 & 0.0145 & -0.0031 \\
    \bottomrule
  \end{tabular}

  \vspace{4pt}
  \begin{minipage}{0.95\linewidth}\footnotesize
  \textbf{Protocol note.} No inference ran. Both columns score the same stored model outputs; only the renderer that produces the perfect-text input, and the scoring rules applied to it, changed. This is the discipline that makes an apparatus repair distinguishable from a model improvement: had these outputs been regenerated, the two causes would have been inseparable. Presence-conditional F$_1$ scores only fields the document actually carries, so it is insensitive to how often fields are absent.
  \end{minipage}
\end{table}


\textbf{The reading-quality proxy was dishonest about blank pages.} A reader that returns
nothing for a page is not a reader that found nothing to report; it is a failure. The proxy
initially treated an empty transcript as a legitimate zero, which allowed a decode failure to
be averaged in with genuine misses instead of being surfaced as a distinct fault.

\textbf{Structural output repair had to be bounded.} A model that emits a \ac{JSON} object with
a missing delimiter has produced a recoverable answer, and refusing to parse it discards a
correct reading. But repair must supply structure only, never values. That boundary is
explicit: delimiters may be inserted, content may not be invented.

\section{Class Four: The Field Specification Was Not Grounded in the Data}
\label{sec:validity-specification}

The fourth class was not anticipated and is, in retrospect, the most instructive.

The specification given to the structuring stage tells it which fields exist and what they tend
to be called on a German invoice. Those printed-label lists are not documentation. \textbf{They
are empirical claims about what invoices print}, and they had been written from domain
knowledge, not measured against the corpus.

Auditing them against all 146 documents found that \textbf{34 of 63 label variants never
occurred in the corpus at all}. Each was a false claim, and each cost specification budget that
would otherwise have gone to a label the corpus actually prints. The failure was not uniform:
the two fields that had been scoring zero on perfect text were precisely the two whose only
German label was the standard's own term, which no invoice in the corpus prints. Meanwhile a
label appearing on 91 of 146 documents, for a field present on 138, was missing from the list
entirely.

Worse, the audit found \textbf{four reference values embedded in the specification as
illustrative examples}. A concrete example value in a field description is, by construction, a
candidate answer: it hands the model the reference for any document that carries it and
inflates that document's score. This is evaluation contamination, and it was introduced by the
same well-meant impulse that produces good documentation.

Two details of how that was found matter more than the finding.

\textbf{The fourth leak was introduced while removing the third.} Replacing a leaked German
payment-method phrase with an English one substituted one reference value for another. Care was
not sufficient; the defect recurred within minutes, in the act of fixing it.

\textbf{It was caught only after the detector was strengthened.} The first detector compared
descriptions against reference values in canonical form. A date written in German day-first
notation does not match its ordered-numeric reference form as a string, so a date-shaped leak
passed the check. Comparing against \emph{printed} variants instead of canonical ones is what
surfaced it.

The conclusion drawn was that vigilance is the wrong instrument for this class of defect. It is
now gated by two mechanisms instead of watched: a self-contained test that rejects
value-shaped content in any field description, and a corpus-backed audit that reports the
occurrence count of every label variant. After repair, no ungrounded variant and no leaked
value remains.

\textbf{One further footgun was found and is recorded, not resolved.} A diagnostic that
constructs the experiment configuration with its defaults, instead of loading the project's
configuration file, silently selects a superseded reading model, because the canonical choice
lives in the configuration, not in the default. The first audit run did exactly that and
therefore measured the wrong pipeline --- a weak reader makes genuinely printed labels look
ungrounded. Both audit tools now load the configuration explicitly. The stale default remains a
live hazard for any future script that constructs the configuration bare, and it is listed in
Chapter~\ref{ch:limitations-future-work}.

\section{A Correction That Was Itself a Defect}
\label{sec:validity-selfinflicted}

One correction in this project broke the symmetry rule of \S\ref{sec:validity-discipline}, and
it is reported here because it is the clearest illustration of why that rule
exists.

The sign-convention fix of \S\ref{sec:validity-comparison} was first implemented on the
predicted side only, on the stated premise that the reference always stores these totals
unsigned. That premise is false. A credit note in the corpus carries a negative discount total
in its reference data --- correctly, because on a credit note the quantity is negative. Folding
the sign on one side only turns a correct negative prediction into a positive one and scores it
as a miss: \textbf{precisely the failure mode the fix existed to remove}, reintroduced by the
fix.

The repair folds both sides, scoped to the three fields whose convention warrants it, and
leaves signed the one total that is legitimately signed. Two regression tests now assert the
symmetry and the exception.

The general lesson is not that this particular premise was wrong. It is that a claim about what
the reference data contains is an empirical claim about a corpus, and stating it in a code
comment does not make it true. It has to be checked against the corpus, in the same way the
field specification did.

\section{What Correcting the Instrument Was Worth}
\label{sec:validity-effect}

All of the following was measured on frozen generations. No model ran, so no part of these
differences is attributable to the system under test.

\textbf{On the perfect-text arm}, Table~\ref{tab:oracle-renderer-correction} gives the effect
of repairing the renderer: every metric improves, the spurious-emission rate falls, and two
fields move from scoring nothing to scoring near the top of the range because they had never
been rendered at all.

\textbf{On the reading arm}, Table~\ref{tab:reader-lineage} shows the same reading model scored
before and after the corrections. The gap between those two rows is instrument, not model:
the same transcripts, the same structurer, the same documents.
Figure~\ref{fig:ruler-correction} plots both arms.

\begin{figure}[htbp]
  \centering
  \includegraphics[width=0.92\textwidth]{ruler-correction}
  \caption[Effect of the scoring corrections]{The effect of repairing the measuring apparatus,
  on both arms, computed from identical archived generations. No model ran between the two
  points, so the entire movement is attributable to the instrument rather than to the system
  under test.}
  \label{fig:ruler-correction}
\end{figure}

\textbf{On the held-out corpus}, the protocol note to Table~\ref{tab:heldout-headline} records
what the answer-key correction was worth, with the true-positive and false-positive counts held
constant by construction.

\textbf{And on a decision}, the effect was qualitative, not numerical. On directly
measured reading quality, the two candidate reading models were separated by three thousandths
under the defective ruler, with the general-purpose model ahead. After the corrections ---
applied to identical stored transcripts --- the ordering inverted, with the gap still under a
point (\S\ref{sec:results-reader}). The correct conclusion from an ordering that flips when
the ruler is fixed is not that the second model is better; it is that \textbf{the instrument
does not separate these two candidates at all}, and a selection made on it would have been
arbitrary. That is what forced the reader decision onto a hand-judged review of every
individual miss, which Chapter~\ref{ch:results} reports. End-to-end score, for completeness,
favoured the transcription model before the corrections and was re-taken only for the selected
reader afterwards; Table~\ref{tab:reader-lineage} records that lineage.

\section{The Near-Miss}
\label{sec:validity-confound}

The most consequential measurement problem in this project was not a defect in the scoring
code. It was a comparison between two numbers that were not comparable.

The adaptation experiments had to run on target-class hardware in full numeric precision,
because they exceeded the floor's memory (\S\ref{sec:method-repro}). The deployed baseline runs
at the floor in 4-bit quantised precision. Comparing the full-precision adapter against the
quantised baseline measures the adapter \emph{and} a precision change at the same time.

Table~\ref{tab:precision-confound} gives both readings. The unmatched comparison shows a
difference small enough to read as no measurable harm. The matched comparison, against a
baseline re-measured in the adapter's own precision, shows a regression roughly twenty times
larger. The two effects --- the damage the adapter does and the gain full precision brings ---
very nearly cancel.

Had the baseline not been re-measured, this thesis would have reported that the adaptation was
harmless. It is not. The requirement to re-measure was written down before the runs, as a
precondition, not a check, which is the only reason the cancellation was caught.

Two things follow. The first is narrow: precision must be matched. The second is general and is
the reason this section exists. \textbf{The confound was not visible in either number.} Both
figures were correctly computed, both were reproducible, and their difference was small and
plausible. Nothing about the output signalled that the comparison was invalid. Only a rule
fixed in advance caught it.

\section{What Remains Broken}
\label{sec:validity-open}

Two items are recorded as open and unrepaired, because reporting only the defects that
were fixed would misrepresent the state of the apparatus.

\textbf{One field regressed under the renderer repair and stays regressed.} The supplier
article number scores worse after the perfect-text renderer was corrected than before. The
repair was still right --- it made the renderer faithful to the reference --- but it exposed a
comparison problem in that field that has not been resolved. It is reported as a regression
instead of smoothed away, and Table~\ref{tab:oracle-renderer-correction} covers the aggregate
in which it sits.

\textbf{The stale configuration default remains.} Described in
\S\ref{sec:validity-specification}: a diagnostic that builds its configuration from defaults
selects a superseded reading model. The two audit tools were fixed; the default was not.

\section{What Generalises}
\label{sec:validity-generalises}

Four practices carried the weight here, and none of them is specific to invoices.

\textbf{Archive generations.} A stored prediction can be re-scored under a corrected rule, which
makes an instrument repair measurable in isolation. Almost every number in this chapter depends
on that one habit.

\textbf{Treat the field specification as data to be measured.} Label lists, alias tables and
example values are empirical claims about a corpus. Measure their occurrence, record the count
next to each entry, and never place a concrete value in a specification --- describe the shape
instead. A value-shaped example is by construction a candidate answer.

\textbf{Distinguish the three causes of a zero.} A field scoring nothing on perfect input has
one of three causes: the value is absent from the input, the value is present but the
specification never names it, or the value is present and named and the comparison rule rejects
it. The middle cause is the one that is routinely forgotten, and it is the one that must never
be handed to a fine-tuning run --- otherwise the adapter is credited with a gain that a one-line
specification change would have produced.

\textbf{Gate what recurs.} Where a defect class recurred under active attention --- as the
contamination class did, within minutes, during its own repair --- vigilance has been
demonstrated to be insufficient. Replace it with a test and an audit that fail loudly.
